\chapter{The finiteness theorem} \label{VIII}

\section[A spectral sequence of biduality]
{A spectral sequence of biduality\protect\sfootnotemark}\label{VIII.1}

\renewcommand{\theequation}{\thesection.\arabic{equation}}

Let us state the result\sfootnotetext{The reader familiar with the
language of Verdier's derived categories will recognize the spectral
sequence associated with a \emph{biduality isomorphism}. \Cf \SGA
6~I.} which we wish to reach:

\begin{proposition} \label{VIII.1.1} Let $A$ be a noetherian ring and let $I$ be an ideal of $A$. Set $X=\Spec(A)$ and $Y=V(I)$. Let $M$ be an $A$-module \emph{of finite type and of finite projective dimension}. Let ${\Fa}=\widetilde{M}$ be the $\OX$-Module associated with $M$.
\begin{enumerate}
\item[\textup{1)}] There exists a spectral sequence:
\[
\H_Y(X, {\Fa}) \Longleftarrow \Ext^p_Y(\Ext^{-q}(M, A), A).
\]

\item[\textup{2)}] There exists a spectral sequence:
\[
\SheafH_Y(X, {\Fa}) \Longleftarrow \SheafExt^p_Y(\SheafExt^{-q}({\Fa}, \OX), \OX).
\]
\end{enumerate}
\end{proposition}

Of course, 2) is deduced from 1) by remarking that, if $M$ and $N$ are two $A$-modules of finite type and if one sets ${{\Fa}}=\widetilde{M}$ and ${{\Ga}}=\widetilde{N}$, one has isomorphisms:
\begin{align*}
\SheafH_Y({\Fa}) &\approx \widetilde{\H_Y(X, {\Fa})}, \\
\SheafExt_Y({\Fa}, {\Ga}) &\approx \widetilde{\Ext_Y({\Fa}, {\Ga}}), \\
\SheafExt_{\OX}({\Fa}, {\Ga}) &\approx \widetilde{\Ext_A(M, N)}.
\end{align*}

Let $\ccat$ be the category of $A$-modules, and $\Ab$ that of abelian groups. Let ${\cF}$ be the functor:
\begin{align*}
\cF: \ccat&\to \Ab \quad\text{defined by} \\
 M &\mto \Gamma_Y(\widetilde{M}).
\end{align*}

One knows from Expos\'e \Ref{II} that there exists an isomorphism of $\partial$-functors:
\[
\H^\ast_Y(X, \widetilde{M})\approx \R^\ast {\cF}(M).
\]
Moreover, let $\Ext^\ast_Y$ be the right derived functors in the second argument of
\[{\cF}\circ \Hom\colon \ccat^\circ \times \ccat \to \Ab.
\]
One knows from Expos\'e \Ref{VI} that one has an isomorphism of $\partial$-functors:
\[
\Ext^\ast_Y(M, N)\simeq \Ext^\ast_Y(\widetilde{M}, \widetilde{N}).
\]
Finally, let us retain the following result from Expos\'e \Ref{VI}: if $C$ is an injective $A$-module and if $N$ is an $A$-module of finite type, the sheaf $\SheafHom(\widetilde{N}, \widetilde{C})\approx\widetilde{\Hom(N, C)}$ is \emph{flasque}, hence $\R^1{\cF}(\Hom(N, C))=0$.

It remains for us to prove the following result:

\begin{lemma} \label{VIII.1.2} Let $A$ be a noetherian ring and let $\ccat$ be the category of $A$-modules. Let ${\cF}\colon \ccat\to\Ab$ be a left exact additive functor such that, for every $A$-module $N$ of finite type and every \emph{injective} $A$-module $C$, one has $\R^1{\cF}(\Hom(N, C))=0$. Let $M$ be an $A$-module of finite type and of finite projective dimension. There exists a spectral sequence:
\[
\R^\ast {\cF}(M)\Longleftarrow \Ext^p_{{\cF}}(\Ext^{-q}(M, A), A), \]
where $\Ext^p_{{\cF}}$ denotes the $p$-th right derived functor of ${\cF}\circ \Hom$.
\end{lemma}

\emph{We shall consider only complexes whose differential is of degree $+1$}. By the hypothesis made on $M$, there exists a \emph{projective} resolution of $M$ of finite length:
\[u\colon L^\boule \to M, \]
where, moreover, the $L^p$ are modules of \emph{finite type} and $L^p=0$ if $p\not\in[-n, 0]$. Let moreover $v\colon M\to I^\boule$ be an injective resolution of $M$. I say that
\begin{equation} \label{eq:VIII.1.1} v\circ u\colon L^\boule\to I^\boule
\end{equation}
is an injective resolution of $L^\boule$. It is appropriate to make precise what this means.

\begin{definition} \label{VIII.1.3} Let $X^\boule$ be a complex of $A$-modules; one calls an injective resolution of $X^\boule$ a homomorphism of complexes:
\[x\colon X^\boule\to CX^\boule, \]
such that $CX^p$ is injective for every $p\in \ZZ$, and such that $x$ induces an isomorphism in homology.
\end{definition}

\begin{proposition} \label{VIII.1.4} Every complex bounded on the left, \ie such that there exists $q\in\ZZ$ with $X^p=0$ for $p<q$, admits an injective resolution. Moreover, if $u\colon X^\boule\to Y^\boule$ is a homomorphism of complexes (bounded on the left) and if $x\colon X^\boule \to CX^\boule$ and $y\colon Y^\boule\to CY^\boule$ are injective resolutions of $X^\boule$ and $Y^\boule$, there exists a homomorphism of complexes:
\[Cu\colon CX^\boule\to CY^\boule, \]
unique up to homotopy, such that the diagram:
\[
\xymatrix{ X^\boule \ar[r]^x\ar[d]_u & CX^\boule\ar[d]^{Cu} \\
Y^\boule \ar[r]^y & CY^\boule \\
}\]
is commutative up to homotopy.
\end{proposition}

The proof is left to the reader\sfootnote{\Cf also H.~Cartan \& S.~Eilenberg, \emph{Homological Algebra}, Princeton Math. Series, vol.~19, Princeton University Press, 1956.}.

Let us recall a notation introduced in Expos\'e \Ref{V}.

\subsection*{Notation}
Let $X^\boule$ and $Y^\boule$ be two complexes. One denotes by $\Hom^\boule(X^\boule, Y^\boule)$ the \emph{simple complex} whose component of degree $n$ is
\[(\Hom^\boule(X^\boule, Y^\boule))^n=\prod_{-p+q=n}\Hom(X^p, Y^q)\]
also denoted $\Hom^n(X^\boule, Y^\boule)$, and whose differential is given by:
\begin{gather*}
\partial_n\colon \Hom^n(X^\boule, Y^\boule)\to \Hom^{n+1}(X^\boule, Y^\boule)\\
\partial_n=d'+(-1)^{n+1}d'',
\end{gather*}
where $d'$ and $d''$ are the differentials (of degree $+1$) induced by those of $X^\boule$ and of $Y^\boule$.

\setcounter{equation}{1} Now let $A^\boule$ be the complex defined by $A^p=0$ if $p\neq 0$ and $A^0=A$. Let
\[a\colon A^\boule \to CA^\boule\]
be an injective resolution of $A^\boule$. Consider the double complex:
\begin{equation} \label{eq:VIII.1.2} Q^{p, q}=\Hom(\Hom(L^{-q}, A), CA^p).
\end{equation}
The first spectral sequence of the bicomplex ${\cF}Q^{\boule\boule}$ will give the conclusion of Lemma \Ref{VIII.1.2}.

Set
\begin{equation} \label{eq:VIII.1.3} {L'}^\boule=\Hom^\boule(L^\boule, A^\boule),
\end{equation}
and
\begin{equation} \label{eq:VIII.1.4} P^\boule=\Hom^\boule({L'}^\boule, CA^\boule).
\end{equation}

One sees easily that $P^\boule$ is the simple complex associated with $Q^{\boule\boule}$. Let us compute the abutment of the spectral sequence \ie the homology of ${\cF}P^\boule$. For this, using the fact that $L^\boule$ is projective of finite type in every dimension, one proves that $L^\boule$ is isomorphic to $\Hom^\boule({L'}^\boule, A^\boule)$. From the homomorphism $a\colon A^\boule\to CA^\boule$, one deduces a homomorphism:
\[b\colon \Hom^\boule({L'}^\boule, A^\boule)\to \Hom^\boule({L'}^\boule, CA^\boule), \]
or again, a homomorphism
\begin{equation} \label{eq:VIII.1.5} c\colon L^\boule\to P^\boule.
\end{equation}
This being said, it is easy to see, using the fact that ${L'}^\boule$ is projective of finite type in every dimension and bounded on the left, that (\Ref{eq:VIII.1.5}) is an injective resolution of $L^\boule$. Using Proposition \Ref{VIII.1.4}., one concludes that $P^\boule$ is homotopically equivalent to $I^\boule$, where $I^\boule$ is the injective resolution of $M$ introduced above (\Ref{eq:VIII.1.1}). One deduces that the \emph{abutment} of the first spectral sequence of the double complex ${\cF}Q^{\boule\boule}$, which is $\H^\ast({\cF}P^\boule)$, \emph{is isomorphic} to $\R^\ast {\cF}(M)$.

The initial term of the first spectral sequence of the bicomplex ${\cF}Q^{\boule\boule}$ is:
\[
\E^{p, q}_2 = {'\!\H}^p({''\!\H}^q({\cF}Q^{\boule\boule})).
\]

For every $p\in\ZZ$, $CA^p$ is injective. By the hypothesis made on ${\cF}$, the functor (restricted to the category of modules of finite type):
\[
N \mto {\cF}\Hom(N, CA^p)\]
is exact. Whence one deduces isomorphisms:
\[
{''\!\H}^q({\cF}\Hom({L'}^\boule, CA^p))\approx {\cF}\Hom(\H^{-q}({L'}^\boule), CA^p).
\]
By the definition of $\Ext^\ast_{{\cF}}$ as derived functor of ${\cF}{\circ}\Hom$, one deduces isomorphisms:
\[
\E^{p, q}_2 \approx \Ext^p_{{\cF}}(\H^q({L'}^\boule), A).
\]
Now ${L'}^\boule=\Hom^\boule(L^\boule, A^\boule)$, where $L^\boule$ is a projective resolution of $M$ whence isomorphisms:
\[
\Ext^{q}(M, A)\approx \H^q({L'}^\boule),
\]
which gives the conclusion.\qed
